\subsection{点测度族}

设 X 和 T 是两个局部紧空间，\(\pi\) 是 T 到 X 中的一个映射，\(g\) 是定义在 T 上的一个有限数值函数 \(\geqslant 0\) ；这两个函数定义了映射 \(t \mapsto \lambda_t = g(t)\varepsilon_{\pi(t)}\) （T 到 X 上测度所成的空间 \(\mathcal{M}(X)\) 中的），使得对每个 \(t \in T\) ，\(\lambda_t\) 或者是一个点测度（Ch. III, §2, No.~4），或者等于 0。如果 \(f\) 是定义在 X 上的一个数值函数 \(\geqslant 0\) ，那么 \(\int^{*} f(x) d\lambda_t(x) = \int^{\bullet} f(x) d\lambda_t(x) = f(\pi(t))g(t)\) （回忆我们的约定：当 \(g(t) = 0\) 且 \(f(\pi(t)) = +\infty\) 时取这个乘积为 0）。定义在 X 上的每个函数（取值于拓扑空间）都是 \(\lambda_t\) -可测的（对每个 \(t \in T\) ）。X 到 Banach 空间 F 中的每个映射 \(\mathbf{f}\) 都是 \(\lambda_t\) -可积的（对所有 \(t \in T\) ），且 \(\int \mathbf{f}(x) d\lambda_t(x) = \mathbf{f}(\pi(t))g(t)\) 。最后，如果 \(f\) 是定义在 X 上的任意数值函数，那么 \(f\) 为 \(\lambda_t\) -可积必须且只须 \(f(\pi(t))g(t)\) 有限，此时 \(\int f(x) d\lambda_t(x) = f(\pi(t))g(t)\) 。

定义 1. --- 设 \(\mu\) 是 T 上一个正测度。称对 \((\pi, g)\) 为 \(\mu\) -适应对，如果下列条件成立：

\(1^{\circ}\) 函数 \(\pi\) 和 \(g\) 都是 \(\mu\) -可测的。

\(2^{\circ}\) 对每个函数 \(f \in \mathcal{K}(X)\) ，映射 \(t \mapsto f(\pi(t))g(t)\) 都是本质 \(\mu\) -可积的。

命题 1. --- 如果对 \((\pi, g)\) 是 \(\mu\) -适应的，那么映射 \(\Lambda: t \mapsto \lambda_t = g(t) \varepsilon_{\pi(t)}\) （T 到 \(\mathcal{M}_+(\mathrm{X})\) 中的）是标量本质 \(\mu\) -可积的、淡 \(\mu\) -可测的且 \(\mu\) -适宜的。反之，如果 \(\Lambda\) 是标量本质 \(\mu\) -可积且淡 \(\mu\) -可测的，那么函数 g 是 \(\mu\) -可测的，且 \(\pi\) 在集 S（由 \(t \in T\) 中使 \(g(t) \neq 0\) 者组成的）上的限制是 \(\mu\) -可测的。

假设对 \((\pi,g)\) 是 \(\mu\) -适应的；对每个函数 \(f \in \mathcal{K}(X)\) ，函数 \(t \mapsto \langle f, \lambda_t \rangle = f(\pi(t))g(t)\) 都是本质 \(\mu\) -可积的。让我们证明 \(t \mapsto \lambda_t\) 是淡 \(\mu\) -可测的。因为，我们首先注意到，如果 \(\pi\) 和 g 连续，那么映射 \(t \mapsto \lambda_t\) 是淡连续的。在一般情形，T 中使 \(\pi\) 和 g 在 K 上的限制都连续的紧子集 K 所成的集是 \(\mu\) -稠密的（Ch. IV, §5, No.~10, 命题 15）；如果 K 是这样一个集，那么 \(t \mapsto \lambda_t\) 在 K 上的限制是淡连续的，由此得出陈述的第一个断言。§3, No.~1 的命题 2 表明 \(\Lambda\) 是 \(\mu\) -适宜的。

反之，假设 \(\Lambda\) 是标量本质 \(\mu\) -可积且淡 \(\mu\) -可测的；那么它是 \(\mu\) -适宜的（§3, No.~1, 命题 2）。由于函数 1 在 X 上是下半连续的，函数 \(t \mapsto \lambda_t(1) = g(t)\) 是 \(\mu\) -可测的（§3, 定义 1）。因此集 S 是可测的（Ch. IV, §5, No.~5, 命题 7）。紧集 K ⊂ S 所成的集 \(\mathfrak{K}\) （这些 K 使 \(g|K\) 连续且 \(\Lambda|K\) 淡连续）在 S 中是 \(\mu\) -稠密的（Ch. IV, §5, No.~10, 命题 15）；如果 K ∈ \(\mathfrak{K}\) ，那么映射 \(t \mapsto \varepsilon_{\pi(t)} = \frac{1}{g(t)} \lambda_t\) 在 K 上的限制是淡连续的，这就蕴涵 \(\pi|K\) 的连续性（Ch. III, §1, No.~9, 命题 13）。由于 \(\mathfrak{K}\) 在 S 中 \(\mu\) -稠密，\(\pi\) 在 S 上的限制是 \(\mu\) -可测的。

我们将利用下面的引理：

引理. --- 设 T 和 X 是两个拓扑空间，\(\pi\) 是 T 到 X 中的一个逆紧连续映射（GT, I, §10, No.~1, 定义 1）。设 g 是定义在 T 上的一个下半连续数值函数。对每个 \(x \in X\) ，令 \(f(x)\) 为函数 \(g(t)\) 在集 \(\overline{\pi}^{-1}(x)\) 中的下确界（下确界等于 \(+\infty\) ，如果 \(\overline{\pi}^{-1}(x) = \varnothing\) ；参见 S, III, §1, No.~9）。则 f 在 X 上是下半连续的。

对每个（有限）实数 \(a\) ，用 \(\mathrm{B}_a\) 表示由 \(x \in \mathrm{X}\) 中使 \(f(x) \leqslant a\) 者组成的集，用 \(\mathrm{A}_a\) 表示由 \(t \in \mathrm{T}\) 中使 \(g(t) \leqslant a\) 者组成的集；问题归结为证明 \(\mathrm{B}_a\) 是闭的（GT, IV, §6, No.~2, 命题 1）。现在，\(\mathrm{A}_a\) 是闭的（同处引证），且逆紧映射 \(\pi\) 是闭的（GT, I, §10, No.~1, 命题 1）；于是我们被归结为证明 \(\pi(\mathrm{A}_a) = \mathrm{B}_a\) 。明显的关系 \(f(\pi(t)) \leqslant g(t)\) （对所有 \(t \in \mathrm{T}\) ）蕴涵 \(\pi(\mathrm{A}_a) \subset \mathrm{B}_a\) 。另一方面，设 \(x \in \mathrm{B}_a\) ；集 \(\overline{\pi}^{-1}(x)\) 是拟紧的（GT, I, §10, No.~2, 定理 1）且非空，因此存在 \(t \in \overline{\pi}^{-1}(x)\) 使 \(g(t) = \inf_{u\in \overline{\pi}^{-1} (x)}g(u) = f(x)\) （GT, IV, §6, No.~2, 定理 3）；从而 \(t\in \mathbf{A}_a\) 且 \(\pi (t) = x\)。

